\usepackage{fontspec}
\usepackage[english, russian]{babel}
\usepackage{graphicx}
\usepackage[a4paper,
            left=3cm,
            right=3cm,
            top=2cm,
            bottom=3cm]{geometry}

\setmainfont{CMU Serif}
\setsansfont{CMU Sans Serif}
\setmonofont{CMU Typewriter Text}

\usepackage{titlesec}

\setcounter{secnumdepth}{4}
\setcounter{tocdepth}{4}

\titleformat{\paragraph}
  {\normalfont\normalsize\bfseries} % Стиль шрифта (обычный, стандартный размер, жирный)
  {\theparagraph}                   % Нумерация
  {1em}                             % Отступ от номера до текста заголовка
  {}                                % Дополнительный код

\titlespacing*{\paragraph}
  {0pt}{3.25ex plus 1ex minus .2ex}{1.5ex plus .2ex}

\newcommand{\subsubsubsection}[1]{\paragraph{#1}}

\usepackage{hyperref}
\hypersetup{
    colorlinks=true,        % цветные ссылки (false - черные рамки)
    linkcolor=black,        % цвет ссылок на разделы
    urlcolor=blue,          % цвет URL-ссылок
    citecolor=blue,         % цвет ссылок на литературу
    pdfborder={0 0 0}       % убрать рамки вокруг ссылок
}

\usepackage[
    backend=biber,
    style=gost-numeric,
    sorting=none
]{biblatex}
\addbibresource{refs.bib}